\appendix
\section{程序与支撑材料}

\subsection{程序文件说明}
\begin{table}[H]
  \centering
  \caption{程序文件与用途}
  \begin{tabularx}{\textwidth}{p{0.28\textwidth}X}
    \toprule
    文件 & 用途 \\
    \midrule
    \texttt{code/example.py} & 程序入口占位示例；定稿时替换为实际提交文件及对应运行说明 \\
    \placeholder{程序文件名} & \placeholder{说明对应问题与运行方式} \\
    \bottomrule
  \end{tabularx}
\end{table}

\lstinputlisting[caption={程序入口示例},label={lst:example}]{code/example.py}

\subsection{支撑材料清单}
\begin{enumerate}
  \item 问题三三次正式测试的加密行为日志，必须保留模拟器导出时的原文件名；
  \item 问题四三次正式测试的加密行为日志，必须保留模拟器导出时的原文件名；
  \item 可复现实验所需的源程序、参数文件与运行说明；
  \item \placeholder{其他支撑材料及其与正文的对应关系}。
\end{enumerate}
